\section{Introduction}
\label{sec:introduction}

The rapid proliferation of shared electric scooters (e-scooters) has fundamentally reshaped urban micromobility across the globe. Since the launch of early dockless operators in 2017, e-scooter sharing services have expanded to over 600 cities worldwide, with annual ridership exceeding 250 million trips in 2023 alone \citep{nacto2024shared}. E-scooters offer a compelling first-and-last-mile transportation solution: they are affordable, require no parking infrastructure, and produce zero tailpipe emissions. However, this rapid expansion has been accompanied by mounting safety concerns. Emergency department visits related to e-scooter injuries have surged in cities across North America, Europe, and Asia, with head trauma, fractures, and soft-tissue injuries accounting for the majority of presentations \citep{trivedi2019injuries, bloom2021standing, badeau2019emergency}. Fatalities, though less frequent, have been documented in multiple jurisdictions and often involve collisions with motor vehicles at speed \citep{yang2020safety}. Pedestrian conflicts on shared sidewalks and bike paths constitute an additional and growing source of concern, particularly for elderly and disabled populations \citep{james2019pedestrians}.

Speed is widely recognized as a fundamental determinant of crash severity across all transportation modes. The power model of the speed-crash relationship, originally established for motor vehicles \citep{elvik2013reexamination}, predicts that even modest increases in operating speed lead to disproportionate increases in injury risk, a relationship that is especially consequential for unprotected road users such as e-scooter riders. Recognizing this, jurisdictions around the world have imposed regulatory speed limits on e-scooters: 25~km/h in South Korea and France, 20~km/h in Germany and the United Kingdom, and varying limits (15--25~km/h) across U.S. cities \citep{ma2021municipal, sexton2023shared}. Yet a critical gap persists between the existence of these speed limits and empirical evidence regarding actual compliance. How often do riders exceed regulatory speed limits? Where and when does speeding concentrate? What rider characteristics and infrastructure features are associated with elevated speeding risk? Answering these questions is essential for evidence-based regulation, but doing so requires large-scale observational data on actual riding speed behavior---data that have, until recently, been largely unavailable.

Prior empirical work on e-scooter speed behavior has been constrained in several important respects. Epidemiological studies document injury patterns and severity but typically lack behavioral data on the rides that produced those injuries \citep{trivedi2019injuries, bloom2021standing}. Survey-based studies capture self-reported attitudes toward speed and risk but cannot measure actual riding behavior \citep{haworth2021comparing}. Naturalistic riding studies using instrumented scooters or on-board GPS loggers provide detailed within-trip speed profiles but are limited to small samples (typically $N < 500$ trips) drawn from a single city \citep{ma2021escooter, tuncer2020notes}. Several studies have leveraged operator-provided GPS trajectory data at a larger scale to analyze spatial patterns of e-scooter usage \citep{jiao2020shared, mckenzie2019spatiotemporal, bai2020dockless}, but these have focused primarily on trip origins, destinations, and route choice rather than on speed behavior. Crucially, no prior study has analyzed within-trip speed profiles---the full sequence of instantaneous speed readings along a ride---at a scale sufficient to characterize city-wide and multi-city patterns of speeding behavior. This limitation prevents researchers and policymakers from understanding the fine-grained speed dynamics that determine crash risk: not merely whether a trip involves speeding at some point, but how much, how often, and in what spatial and temporal contexts.

In this paper, we address these gaps by presenting the first large-scale empirical characterization of e-scooter riding speed behavior using operator telemetry data from Swing, a major shared e-scooter operator in South Korea. Our primary dataset comprises 2.83~million GPS trajectories collected during May 2023, each containing a full within-trip speed vector recorded at approximately 10-second intervals. We supplement this with 44.1~million trip records spanning 24 months (January 2022 through December 2023) for longitudinal experience analysis, and 20~million trip-level records with user demographics, vehicle characteristics, and billing information. After rigorous quality filtering---removing GPS errors, implausible speeds, and trips with insufficient data points---we retain 2.78~million valid cross-sectional trips (98.2\% retention) covering 382,328 unique users, 81,525 unique vehicles, and 52 cities across 17 provinces in South Korea. To our knowledge, this is the largest dataset of e-scooter trips with per-point speed profiles ever analyzed in the academic literature, exceeding the scale of prior GPS-based e-scooter studies by two to three orders of magnitude.

We make nine contributions to the e-scooter safety literature:

\begin{enumerate}
    \item \textbf{Unprecedented empirical scale.} We provide the first city-scale characterization of e-scooter speed behavior using 2.78~million trajectories with full within-trip speed profiles from 52 cities. We find that 29.6\% of all trips contain at least one speed reading exceeding the Korean regulatory limit of 25~km/h, while 60.9\% of users never exceed this threshold.

    \item \textbf{Comprehensive speed indicator framework.} We define and compute a multi-level suite of safety-relevant speed indicators---at the trip, user, road segment, and spatial grid levels---including speeding exposure rate, speed variability, harsh acceleration frequency, and cruise behavior metrics. This framework provides a standardized, reproducible methodology for speed-based micromobility safety assessment.

    \item \textbf{Spatial analysis of speeding patterns.} Using H3 hexagonal grid aggregation and spatial statistics (Getis-Ord $G_i^*$ and Moran's $I$), we identify statistically significant spatial clusters of elevated speeding across five major Korean cities. Global Moran's $I$ values range from 0.16 to 0.71 (all $p < 0.001$), confirming strong spatial autocorrelation of speeding behavior. We overlay hotspot maps with road infrastructure attributes to characterize the infrastructure profiles associated with high-speeding zones.

    \item \textbf{Multi-factor statistical modeling.} We estimate a suite of regression models---logistic with generalized estimating equations (GEE), mixed-effects linear, and a two-part hurdle model with beta regression---to disentangle the contributions of rider demographics (age), vehicle operating mode, temporal context, and road infrastructure to speeding propensity and intensity. The vehicle operating mode emerges as the dominant predictor: the turbo (TUB) mode, which removes the software-based speed governor, increases the odds of speeding by a factor of 121 relative to the standard (STD) mode (OR = 121, $p < 0.001$).

    \item \textbf{Natural experiment on speed governors.} We exploit the within-user variation between TUB (ungoverned) and ECO (governed) operating modes as a natural experiment. Among 12,046 users who used both modes, the ECO speed governor reduces mean speed by 24\%, maximum speed by 31\%, and the speeding rate from 7.4\% to 0.38\% (all $p < 0.001$, Cohen's $d = 1.0$--$2.2$). This provides direct causal evidence on the effectiveness of hardware-level speed limiting as a safety intervention.

    \item \textbf{Data-driven rider typology.} We identify four distinct rider speed behavior typologies via Gaussian mixture modeling: Safe Riders (40\%), Moderate-Risk Riders (26\%), Habitual Speeders (26\%), and Stop-and-Go Riders (8\%). Multinomial logistic regression reveals that TUB mode usage, younger age, and specific geographic locations strongly predict membership in higher-risk classes, enabling targeted safety interventions.

    \item \textbf{Longitudinal experience--speeding analysis.} Using 44.1~million trips across 24 months, we demonstrate that speeding increases monotonically with rider experience (from 3.8\% on the first trip to 28.0\% after 500+ trips), driven primarily by progressive adoption of the ungoverned TUB mode rather than within-mode behavior changes. This finding supports graduated mode access policies for new riders.

    \item \textbf{Trip geometry, road composition, and model interpretability.} We quantify the non-linear relationship between trip distance and speeding (pseudo-$R^2 = 0.387$), identify differential effects of road class composition on speeding risk, and reveal a curvature--speeding risk paradox (lower prevalence but higher risk on curvy roads). A SHAP analysis based on gradient-boosted trees (AUC $= 0.905$) validates the parametric regression findings and confirms the dominance of operating mode.

    \item \textbf{Quasi-experimental causal evidence.} We exploit Swing's December 2023 system-wide ban of TUB mode as a natural experiment, applying a difference-in-differences design across 53 cities to estimate the population-level causal effect of speed governor policy. Each 10 percentage point reduction in TUB availability produces a 5.0 percentage point reduction in city-level speeding (95\% CI: [4.1, 7.4]~pp), with a null placebo test ($p = 0.56$) supporting causal validity. This provides the first quasi-experimental evidence on the causal effect of e-scooter speed governors at scale.
\end{enumerate}

Our findings carry direct implications: mandatory speed governors could reduce speeding $\sim$20-fold; spatial clustering supports targeted geofencing; rider typologies enable graduated enforcement; and the experience--speeding escalation, mediated by mode adoption, motivates graduated mode access policies. More broadly, we demonstrate that operator telemetry data can provide the evidence base cities need for calibrating e-scooter speed regulations.
